Os dados brutos que compôem um \textit{trace} são similares a um \textit{snapshot}, isto é, uma representação pontual do estado de um sistema. Dessa forma, para a facilitar a ingestão de dados, livre de \textit{constraints} relacionais, optou-se pelo banco MongoDB para armazenamento e consulta dos \textit{logs} de atividade. Para o caso do ESPEON, foi criado o banco TRACES, com os dados sendo estocados na \textit{collection} \textit{logs}. Os registros são estruturados conforme a Figura~\ref{fig:logs_example}, a seguir. 

\begin{figure}[H]
    \centering
    \begin{minipage}{0.9\textwidth}
        \centering
        \includegraphics[width=.9\textwidth]{espeon/exemplo.log.png}
        \caption{Exemplo de \textit{log}}
        \label{fig:logs_example}
    \end{minipage}\hfill
\end{figure}

Para demais componentes da solução, dada a necessidade de cruzar referências e estabelecer relacionamentos entre entidades, foi escolhido um banco relacional: PostgreSQL. Sua escolha se deu pelas facilidades que o banco possui e pelo suporte que o serviço de hospedagem escolhido para implantação o confere. O banco é composto por seis tabelas distintas, cujas especificidades estão dispostas a seguir.

\pagebreak
\subsection{Tabela Users}
A tabela \textbf{Users} armazena os registros de usuários do sistema; usada para identificação e caracterização de cada indivíduo. É nessa tabela que a função do usuário na plataforma é definida, podendo ser \textit{student} ou \textit{educator}, aluno e docente, respectivamente. 
\begin{figure}[H]
    \centering
    \begin{minipage}{0.7\textwidth}
        \centering
        \includegraphics[width=.9\textwidth]{espeon/table.users.png}
        \caption{Tabela Usuários}
        \label{fig:pg_users}
    \end{minipage}\hfill
\end{figure}
\pagebreak

\subsection{Tablea Login}
A tabela \textbf{Login} tem como finalidade associar um usuário a um \textit{token} de autenticação. Nela são registrados o \textit{token} gerado e a respectiva data de criação, permitindo validar a autenticidade e a vigência das sessões de acesso. Esse mecanismo assegura controle sobre o ciclo de vida de autenticações, contribuindo para a segurança da solução.
\begin{figure}[H]
    \centering
    \begin{minipage}{0.7\textwidth}
        \centering
        \includegraphics[width=.9\textwidth]{espeon/table.login.png}
        \caption{Tabela Login}
        \label{fig:pg_login}
    \end{minipage}\hfill
\end{figure}
\pagebreak

\subsection{Tablea Subject}
A tabela \textbf{Subject} reúne os registros de disciplinas cadastradas no sistema. Cada disciplina é identificada por seu nome e \textit{id}, servindo como entidade base para o vínculo com aulas e relatórios.
\begin{figure}[H]
    \centering
    \begin{minipage}{0.7\textwidth}
        \centering
        \includegraphics[width=.9\textwidth]{espeon/table.subject.png}
        \caption{Tabela Subject}
        \label{fig:pg_subject}
    \end{minipage}\hfill
\end{figure}
\pagebreak

\subsection{Tabela Lecture}
A tabela \textbf{Lecture} registra cada aula individual ministrada dentro da plataforma. Cada entrada associa a aula a uma disciplina, a um usuário docente responsável e ao período de realização, especificado pelos campos de início e fim. Essa tabela é fundamental para relacionar eventos de monitoramento e para consolidar informações posteriormente utilizadas na geração dos relatórios de aula.
\begin{figure}[H]
    \centering
    \begin{minipage}{0.7\textwidth}
        \centering
        \includegraphics[width=.9\textwidth]{espeon/table.lecture.png}
        \caption{Tabela Lecture}
        \label{fig:pg_lecture}
    \end{minipage}\hfill
\end{figure}
\pagebreak

\subsection{Tabela Report}
A tabela \textit{Report} armazena as informações consolidadas de cada relatório gerado pelo sistema. Um relatório agrega métricas de engajamento, atenção e participação capturadas durante uma aula e processadas pelo mecanismo de análise de dados. Essa tabela funciona como o repositório central dos resultados analíticos, mantendo vínculo direto com a aula correspondente e servindo como base para consultas e visualizações posteriores.
\begin{figure}[H]
    \centering
    \begin{minipage}{0.9\textwidth}
        \centering
        \includegraphics[width=.9\textwidth]{espeon/table.report.png}
        \caption{Tabela Report}
        \label{fig:pg_report}
    \end{minipage}\hfill
\end{figure}
\pagebreak

% \subsection{Report Student}
% A tabela \textbf{Report Student} estabelece a relação entre um relatório e os estudantes participantes da aula correspondente. Cada registro vincula um relatório a um aluno específico, permitindo armazenar métricas individualizadas, como tempo de visualização, períodos de inatividade e indicadores de engajamento. Essa tabela possibilita análises detalhadas por aluno, fundamentais para acompanhar desempenho e identificar comportamentos relevantes.
% \begin{figure}[H]
%     \centering
%     \begin{minipage}{0.7\textwidth}
%         \centering
%         \includegraphics[width=.9\textwidth]{espeon/table.report_student.png}
%         \caption{Tabela Report Student}
%         \label{fig:pg_rpt_stdt}
%     \end{minipage}\hfill
% \end{figure}
